\chapter{交付与部署}
\label{ch:22}

Magel 程序通过 VM 执行。交付时，可以保留源码和编译器，也可以预先导出字节码并随宿主部署。

\section{直接运行源码}

\par
\medskip
\begin{lstlisting}[style=sh]
magel run app.mgl
\end{lstlisting}
\par
\medskip

标准库内嵌在编译器中。实际使用的动态库、数据库、网络服务和外设仍需要部署环境提供。

\section{导出字节码}

\par
\medskip
\begin{lstlisting}[style=sh]
magel svmb app.mgl out build/app
\end{lstlisting}
\par
\medskip

程序产物包括 \texttt{.\allowbreak{}svmb}、\texttt{.\allowbreak{}data.\allowbreak{}bin} 与 \texttt{.\allowbreak{}layout}。宿主校验字节码格式和特征位，加载数据，注册原生函数，然后通过 \texttt{stackvm:\allowbreak{}:\allowbreak{}Machine} 执行。

四元式可通过 \texttt{magel emit app.\allowbreak{}mgl out build/\allowbreak{}app} 导出，用于检查或独立 VM 工具装配。四元式和字节码都不是宿主机器码。

\section{嵌入宿主}

桌面应用和嵌入式 Linux 宿主可以编译共享 VM，内嵌字节码及数据。宿主程序本身由 C/C++ 工具链构建，Magel 逻辑仍通过 VM 执行。

宿主负责程序、数据区、注册表、堆与输出接口的生命周期，检查执行退出状态、错误和活动资源。共享 VM 核心只有 \texttt{library/\allowbreak{}stackvm/\allowbreak{}} 一份，宿主按目标架构构建胶水和系统依赖。

\section{预算、错误与恢复}

VM 默认使用最高可用执行优化。桌面可配置步数与内存预算；嵌入式 Linux 按实际服务与内存容量设置预算。

看门狗支持可恢复终止及宿主重启；重启前必须清理 VM 所有权状态，不能把旧地址或闭包当成新程序中的有效对象。

\section{符号压缩与程序更新}

符号压缩降低符号表体积和可读性，不构成加密、防逆向或完整沙箱。交付时仍应按实际宿主权限控制文件、网络及外设能力。

程序更新替换字节码映像，不自动迁移运行中的栈、闭包和全局状态。先停止旧实例，清理资源，再装入匹配的数据和新程序。

\section{命令边界}

工具选项见\href{https://github.com/rabbitc999/Magel/blob/master/doc/09-toolchain.md}{工具链}，宿主代码见\href{https://github.com/rabbitc999/Magel/blob/master/doc/08-host-integration.md}{宿主集成}，目标宿主部署见\hyperref[ch:25]{第 25 章}。

下一章通过\hyperref[ch:23]{命令行待办项目}串起语言和标准库。
